\section{Open Questions}

The following five questions were left genuinely open by the replication. Each
is paired with a concrete next-step probe that a follow-up experiment could
run without needing hardware access or paid compute.

\begin{enumerate}
  \item \textbf{qDRIFT vs.\ higher-order Trotter.} Does qDRIFT retain its
    $L$-independent gate-count advantage against Trotter-2 and Trotter-4
    (higher-order Suzuki formulas), or only against first-order Trotter?

    \emph{Basis.} The replication compared qDRIFT only against Trotter-1,
    which is the weakest deterministic baseline. Higher-order Suzuki formulas
    have better $t$ scaling and are known to be competitive with qDRIFT in
    low-commutator regimes. A quantitative crossover map is not fully
    established in the reproduced experiments.

    \emph{Next steps.} Extend the 4-qubit dense simulator to include
    Suzuki-Trotter of orders 2 and 4. Sweep $(L, \lambda, t)$ with $\lambda$
    held fixed and measure the crossover $N$ at which qDRIFT beats each order
    for random Pauli Hamiltonians, then repeat for a structured Hamiltonian
    (e.g.\ 1D Heisenberg or 4-site Fermi--Hubbard via Jordan--Wigner). Produce
    a phase diagram in $(L, t)$ space.

  \item \textbf{Time-dependent Hamiltonians.} How does qDRIFT scale for
    time-dependent Hamiltonians $H(t)$, which the paper's static analysis
    does not directly cover?

    \emph{Basis.} The paper analyzes time-independent $H$. Many practical
    targets (adiabatic state prep, driven systems, dynamical decoupling) are
    time-dependent. Follow-up work has proposed time-dependent generalizations,
    but a clean direct extension with matching error bounds is nontrivial.

    \emph{Next steps.} Implement a naive time-dependent qDRIFT by
    (a) partitioning $[0, t]$ into $K$ windows, (b) resampling term
    probabilities $p_j(t_k) = h_j(t_k) / \lambda(t_k)$ per window, and
    (c) measuring channel error vs total sample count for a driven 4-qubit
    Hamiltonian (e.g.\ transverse-field Ising with time-dependent field).
    Compare against a piecewise-Trotter baseline; look for whether
    $\bar{\lambda} := (1/t) \int_0^t \lambda(s)\,ds$ controls the bound.

  \item \textbf{Composition with importance-sampled multiproduct formulas.}
    Can qDRIFT be composed with importance-sampled multiproduct formulas to
    reduce the $\lambda^2$ factor?

    \emph{Basis.} qDRIFT's cost is quadratic in $\lambda$, which dominates
    for chemistry Hamiltonians where $\lambda$ scales with molecular size.
    Multiproduct formulas and importance sampling of coefficients have been
    explored separately; a hybrid that inherits qDRIFT's $L$-independence
    while getting sub-quadratic $\lambda$ scaling is not obviously ruled out.

    \emph{Next steps.} Prototype a two-stage compiler:
    (1) group Hamiltonian terms into $O(\log 1/\varepsilon)$ blocks by
    coefficient magnitude, (2) apply a multiproduct Trotter formula across
    blocks, (3) sample within each block a la qDRIFT. Measure empirical
    $\lambda$-scaling on random Hamiltonians with a fat-tailed coefficient
    distribution and check whether the effective exponent drops below 2.

  \item \textbf{Noise robustness.} How robust is qDRIFT's advantage under
    realistic gate noise (depolarizing, coherent overrotation, cross-talk),
    where each additional sampled gate injects noise as well as signal?

    \emph{Basis.} The paper's analysis is fully coherent / noiseless. Since
    qDRIFT typically needs more gates than an optimized high-order Trotter to
    reach the same channel error, and noise cost is per-gate, the noiseless
    advantage may be eroded or reversed on near-term hardware.

    \emph{Next steps.} Extend the 4-qubit simulator to a noisy channel model:
    after each sampled gate, apply a single-qubit depolarizing channel of
    strength $p \in \{10^{-4}, 10^{-3}, 10^{-2}\}$ and a small coherent
    overrotation of angle $\delta$. Compare the resulting total channel error
    against Trotter-2 with the same per-gate noise. Identify the noise
    threshold at which qDRIFT loses its $L$-independent advantage for
    $L \in \{24, 60, 200\}$.

  \item \textbf{Adaptive-sampling variants.} Do adaptive-sampling variants of
    qDRIFT (where sample probabilities are updated online based on measured
    or estimated residual error) achieve provable improvements over the
    static $p_j = h_j / \lambda$ sampling, and if so under what structural
    assumptions on $H$?

    \emph{Basis.} qDRIFT samples i.i.d.\ from a static distribution. Adaptive
    or low-discrepancy sampling schemes (quasi-Monte Carlo, control variates,
    stratification) are natural candidates for variance reduction. Whether
    they yield a formal asymptotic improvement over
    $2 \lambda^2 t^2 / \varepsilon$, or only constant-factor gains, is open.

    \emph{Next steps.} Implement three variants:
    (a) stratified sampling guaranteeing term $j$ appears at least
    $\lfloor N p_j \rfloor$ times,
    (b) a control-variate scheme that subtracts a cheap Trotter-1 evolution
    and applies qDRIFT to the residual, and
    (c) an online adaptive scheme reweighting $p_j$ based on estimated
    per-term commutator contributions.
    Measure channel error vs $N$ on the same 4-qubit benchmarks; any variant
    showing a scaling exponent change (not merely a prefactor drop) is a
    candidate for a theoretical result.
\end{enumerate}
